\documentclass[11pt,a4paper]{article}

\usepackage[T1]{fontenc}
\usepackage[utf8]{inputenc}
\usepackage[spanish,es-noshorthands]{babel}
\usepackage{lmodern}
\usepackage{amsmath,amssymb,amsthm}
\usepackage[margin=2.7cm]{geometry}
\usepackage{booktabs}
\usepackage{microtype}
\usepackage[colorlinks=true,linkcolor=black,citecolor=black,urlcolor=blue]{hyperref}

\theoremstyle{plain}
\newtheorem{teorema}{Teorema}
\newtheorem{corolario}[teorema]{Corolario}
\newtheorem{proposicion}[teorema]{Proposición}
\theoremstyle{definition}
\newtheorem{problema}{Problema abierto}
\theoremstyle{remark}
\newtheorem*{obs}{Observación}

% babel-spanish redefine \% con un chequeo de \lastskip que rompe en modo
% matematico ("Incompatible glue units"); \pc lo evita.
\newcommand{\pc}{\ifmmode\text{\,\%}\else\,\%\fi}

\newcommand{\MER}{\mathrm{MER}}
\newcommand{\dd}{\,\mathrm{d}}
\renewcommand{\qedsymbol}{$\blacksquare$}

\title{\vspace{-1.5cm}\textbf{Un problema inverso para el ascenso bifásico\\ de magma por un conducto volcánico}\\[0.3em]
\large Nota de formalización}
\author{Ignacio Gómez Gómez}
\date{\today}

\begin{document}
\maketitle
\vspace{-0.8cm}

\begin{abstract}\noindent
El modelo directo de la tesis es el de conducto \textbf{bifásico} ya
construido en este repositorio (dos velocidades, arrastre, exsolución,
cristalización). El problema inverso consiste en recuperar, a partir de
observables de superficie, parámetros que ese modelo toma como dato: geometría
del conducto, sobrepresión de cámara, contenido de volátiles y burbujas en
profundidad. Para \emph{demostrar} que esa recuperación no es única se usa un
recorte monofásico del mismo sistema, no un modelo distinto. La pregunta
central de la tesis es si esa no-unicidad \textbf{sobrevive} en el modelo
bifásico.
\end{abstract}

\section{El modelo de la tesis es el bifásico}

El operador directo ya existe: el código \texttt{RIconduitex5\_5.py} (y sus
versiones transiente y 2D) resuelve el ascenso de magma como un flujo de
\textbf{dos fases} ---fundido $m$ y gas $g$--- con velocidades distintas, según
el marco de Slezin (2003) y Kozono \& Koyaguchi (2009). El estado en cada
profundidad $z$ es
\begin{equation}
\mathbf{y}(z)=\bigl(P,\;\phi,\;N,\;\xi,\;u_m,\;u_g\bigr),
\end{equation}
presión, fracción volumétrica de gas, densidad numérica de burbujas, fracción
de cristales y las dos velocidades. El sistema es un DAE de seis ecuaciones:
conservación de masa de cada fase, dos balances de momentum (con fricción
pared--fundido $F_{mw}$ y arrastre gas--fundido $F_{mg}$), cinética de
cristalización y evolución de $N$ por coalescencia. La solubilidad es
$c_s(P)=C_1 P^\beta$; el agua disuelta $c_d$ y el gas exsuelto se calculan
comparando el contenido total $c_0$ con esa capacidad.

Los parámetros que el modelo \emph{recibe} y el observable que \emph{entrega}
definen el operador directo
\begin{equation}
\mathcal{F}:\ \Theta\longrightarrow\mathcal{D},\qquad
\theta=\bigl(R(\cdot),\,P_{cam},\,c_0,\,T,\,\xi_0,\,\phi_{crit},\,\ldots\bigr)
\longmapsto d=\MER,
\end{equation}
junto, si se desea, con perfiles $\mathbf{y}(z)$ o con la serie $\MER(t)$ en
la versión transiente. $\Theta$ es de dimensión infinita en cuanto contiene
la función $R(\cdot)$. Caso base: Calbuco 2015
(\texttt{calbuco2015d.py}).

\textbf{Ese es el modelo que se invierte.} No se reemplaza por uno de una sola
velocidad. Las extensiones transiente y 2D del mismo repositorio se mantienen
como variantes del mismo operador, no como un proyecto aparte. La formulación
completa ---masa, los dos momenta, solubilidad, cristales, coalescencia,
regímenes y el DAE--- está escrita en
\texttt{tesis/matematica\_modelo\_bifasico.pdf}.

\section{El problema inverso}

Dado un dato $d^{obs}=\mathcal{F}(\theta^\dagger)+\varepsilon$, las tres
preguntas clásicas son:

\begin{enumerate}
\item[(P1)] \textbf{Unicidad.} ¿Es $\mathcal{F}$ inyectivo? Si no, caracterizar
  $\mathcal{F}^{-1}(d)$.
\item[(P2)] \textbf{Estabilidad.} ¿Es $\mathcal{F}^{-1}$ continuo? ¿Con qué
  módulo?
\item[(P3)] \textbf{Reconstrucción.} Algoritmo, regularización, cuantificación
  de incertidumbre.
\end{enumerate}

Las incógnitas de interés geológico son, entre otras, la presión en
profundidad $P_{cam}$, el contenido de burbujas $\phi(H)$, el agua total $c_0$
y la geometría $R(\cdot)$. La respuesta a (P1) es negativa ya en un recorte
del modelo, y se puede demostrar. La tesis pregunta si sigue siéndolo en el
bifásico.

\section{Un recorte, no un reemplazo}

Para \emph{demostrar} (P1) hace falta un operador en el que se pueda separar
variables. Se obtiene recortando el bifásico, no cambiándolo:

\begin{center}
\begin{tabular}{@{}lll@{}}
\toprule
 & \textbf{Modelo de la tesis (bifásico)} & \textbf{Recorte para el teorema} \\
\midrule
velocidades & $u_m\neq u_g$ (arrastre $F_{mg}$) & $u_m=u_g$ (mezcla) \\
cristales & cinética $\mathrm{d}\xi/\mathrm{d}z$ & $\xi$ fijo \\
exsolución & $c_d$ puede ir atrasada & equilibrio $c_d=\min(c_0,C_1P^\beta)$ \\
dominio & conducto completo, con fragmentación & tramo viscoso hasta $z_f$ \\
geometría & cilindro o dique (hoy) & $R(z)$ arbitraria \\
constitutivas & Bottinga--Weill, Giordano, \texttt{fvrel} & las mismas \\
\bottomrule
\end{tabular}
\end{center}

\vspace{0.4em}
El recorte comparte las constitutivas y los parámetros de Calbuco. Su único
rol es permitir un teorema; el operador que se quiere invertir sigue siendo
$\mathcal{F}$ bifásico. En el recorte, $\rho=\rho(P)$ y $\mu=\mu(P)$, el
caudal $q=\rho A u$ es el $\MER$, y el momentum de lubricación da
\begin{equation}\label{eq:momentum}
\frac{dP}{dz}=-\rho(P)\,g-\frac{8\,\mu(P)\,\MER}{\pi\,\rho(P)\,R(z)^4},
\end{equation}
con $P(H)=P_{cam}$ y $P(z_f)=P_{frag}$.

\section{Resultado del recorte: degeneración geométrica}

En el régimen dominado por fricción (el término $\rho g$ despreciable frente
al viscoso) se define
\begin{equation}
J_4[R]:=\int_H^{z_f}R(z)^{-4}\dd z.
\end{equation}

\begin{teorema}[Degeneración]\label{th:J4}
En ese régimen,
\begin{equation}\label{eq:invariante}
\MER\cdot J_4[R]=\frac{\pi}{8}\,\Phi(P_{cam},P_{frag}),
\qquad
\Phi:=\int_{P_{frag}}^{P_{cam}}\frac{\rho(P)}{\mu(P)}\dd P .
\end{equation}
El $\MER$ depende de $R(\cdot)$ \textbf{únicamente a través del escalar}
$J_4[R]$.
\end{teorema}

\begin{proof}
Despreciando $\rho g$ en \eqref{eq:momentum} y separando variables,
\[
\frac{\rho(P)}{\mu(P)}\dd P=-\frac{8\,\MER}{\pi}\,\frac{\dd z}{R(z)^4}.
\]
Integrando de $z=H$ a $z=z_f$ se obtiene \eqref{eq:invariante}. El lado
izquierdo no depende de $R(\cdot)$.
\end{proof}

\begin{corolario}[No inyectividad]
$\mathcal{F}_{\mathrm{recorte}}$ no es inyectivo en $R(\cdot)$. El conjunto
de nivel $\{R:J_4[R]=c\}$ es de dimensión infinita, y todos sus elementos
producen exactamente el mismo dato.
\end{corolario}

\begin{corolario}[Espacio nulo]
$DJ_4[R]\,\delta R=-4\int R^{-5}\,\delta R\dd z$, luego
\begin{equation}
\mathcal{N}_R=\Big\{\delta R\in C([H,z_f])\ :\ \int_H^{z_f}R(z)^{-5}\,\delta R(z)\dd z=0\Big\}
\end{equation}
es un subespacio cerrado de \textbf{codimensión~1}.
\end{corolario}

\begin{obs}
De las infinitas direcciones en que se puede perturbar la geometría, todas
menos una son invisibles para la tasa de descarga. Eso es una afirmación
sobre el recorte. En el bifásico el arrastre $F_{mg}$ y el deslizamiento
$u_g-u_m$ \emph{podrían} romperla: ésa es la pregunta de la tesis, no un
resultado ya cerrado.
\end{obs}

El resultado no es un artefacto de Poiseuille. Basta que la fricción
factorice como $dP/dz=-F(P)\,f(\MER)\,R^{-n}$: entonces el dato depende de
$R$ sólo a través de $J_n[R]=\int R^{-n}\dd z$, con $n=4$ (laminar),
$n=19/4$ (Blasius) o $n=5$ (Darcy--Weisbach).

\section{Qué aportan (y qué no) los datos temporales}

Con una cámara que se drena, el dato pasa a ser la curva
$P_{cam}\mapsto\MER_{ss}(P_{cam})$. En el recorte eso \textbf{no} informa más
sobre $R(\cdot)$: cada punto de la curva devuelve el mismo $J_4[R]$
(Proposición~\ref{prop:tiempo} más abajo). Sí separa \emph{escalares} que
entran en $\Phi$, por ejemplo el par $(R,\Delta P)$ cuando $R$ es constante.

\begin{proposicion}\label{prop:tiempo}
Bajo las hipótesis del Teorema~\ref{th:J4}, conocer
$P_{cam}\mapsto\MER_{ss}(P_{cam})$ en un intervalo determina $J_4[R]$ y nada
más de $R(\cdot)$.
\end{proposicion}

\begin{proof}
De \eqref{eq:invariante},
$J_4[R]=\tfrac{\pi}{8}\,\Phi(P_{cam},P_{frag})/\MER_{ss}(P_{cam})$ en cada
punto; $\Phi$ es conocida. Todos los puntos dan el mismo escalar.
\end{proof}

Numéricamente, en el recorte, agregar la serie de tiempo reduce el intervalo
de confianza del radio constante de $[8,19]$~m a $[13,18]$~m, y deja la
sobrepresión con sólo una cota superior. Separar dos números y recuperar una
función son problemas distintos.

\section{Verificación numérica del recorte}

Con el recorte, constitutivas del repositorio y parámetros de Calbuco 2015:

\begin{enumerate}
\item Sin flotabilidad, $\MER\cdot J_4$ es constante a $<10^{-6}$ relativo
  entre cinco geometrías, y \eqref{eq:invariante} reproduce el $\MER$
  numérico con error $+0.000\pc$.
\item \emph{Con} flotabilidad el invariante ya no es $J_4$, pero la
  no-unicidad persiste. Ajustando un factor de escala por forma:
\end{enumerate}

\begin{center}
\small
\begin{tabular}{@{}lrrrr@{}}
\toprule
\textbf{Geometría} & $R(H)$ [m] & $R(z_f)$ [m] & \textbf{Volumen} [m$^3$] & \textbf{error en} $\MER$ \\
\midrule
cilindro & 16.0 & 16.0 & $4.83\times10^6$ & $+0.00000\pc$ \\
ensanche basal $\times3$ & 44.4 & 14.8 & $1.29\times10^7$ \ ($+166\pc$) & $-0.00001\pc$ \\
ensanche basal $\times4$ & 51.4 & 13.6 & $2.01\times10^7$ \ ($+317\pc$) & $-0.00002\pc$ \\
ensanchamiento superior $\times2$ & 11.2 & 22.4 & $4.74\times10^6$ & $+0.00000\pc$ \\
dos tramos $\times2$ & 10.7 & 21.5 & $5.38\times10^6$ & $+0.00001\pc$ \\
constricción profunda $-35\pc$ & 16.8 & 16.9 & $4.80\times10^6$ & $+0.00001\pc$ \\
\bottomrule
\end{tabular}
\end{center}

Seis conductos con volúmenes que difieren en $+317\pc$ dan el mismo caudal
con error $\sim10^{-7}$. El $90\pc$ de la resistencia está en los últimos
$15$~m: el conducto profundo es invisible para el $\MER$. \textbf{Esto está
probado en el recorte.} Falta verlo en el bifásico.

\section{Preguntas abiertas (contenido de la tesis)}

\begin{problema}[¿Sobrevive la degeneración en el bifásico?]
Con $u_m\neq u_g$, arrastre $F_{mg}$, permeabilidad y cristalización
cinética, $\mathcal{F}$ ¿sigue siendo no inyectivo en $R(\cdot)$? El
deslizamiento introduce una escala (el radio de burbuja) que el recorte no
tiene, y podría restaurar unicidad o dejar una degeneración de codimensión
distinta. Se busca o bien un contraejemplo numérico en
\texttt{RIconduitex5\_5.py}, o bien un resultado de unicidad.
\end{problema}

\begin{problema}[Grado de mal planteamiento]
Decaimiento de los valores singulares de $D\mathcal{F}[\theta^\dagger]$
restringido a perturbaciones de $R(\cdot)$, primero en el recorte y después
en el bifásico. Eso fija la tasa de convergencia bajo condiciones de fuente.
\end{problema}

\begin{problema}[Unicidad condicional]
¿Qué datos adicionales restauran la unicidad en el bifásico? Candidatos: el
perfil $\phi(z)$ o las texturas de piroclastos (que ``ven'' $u_g-u_m$), la
profundidad de fragmentación, deformación superficial, espectro sísmico del
conducto. Se conjetura un resultado negativo para el espectro por analogía
con Webster y el teorema de dos espectros de Borg (1946).
\end{problema}

\section*{Nota sobre el estado del arte}

La no-unicidad de \emph{parámetros escalares} ya está documentada: Anderson
\& Segall (2013) señalan que de sus datos sólo se determina $R^4/8\eta$, no
$R$ y $\eta$ por separado. Lo que no he visto hecho es (i)~tratar $R(\cdot)$
como incógnita funcional, (ii)~caracterizar el espacio nulo, (iii)~preguntar
si el resultado sobrevive al deslizamiento entre fases, que es precisamente
la física que el modelo de este repositorio ya resuelve.

\begin{thebibliography}{9}\small
\bibitem{anderson2013} Anderson, K. \& Segall, P. (2013).
  \emph{Bayesian inversion of data from effusive volcanic eruptions using
  physics-based models: Application to Mount St.\ Helens 2004--2008}.
  JGR 118:2017--2037. \href{https://doi.org/10.1002/jgrb.50169}{10.1002/jgrb.50169}
\bibitem{castruccio2016} Castruccio, A.\ \emph{et al.} (2016).
  \emph{Eruptive parameters and dynamics of the April 2015 sub-Plinian eruptions
  of Calbuco volcano}. Bull.\ Volcanol.\ 78(9).
  \href{https://doi.org/10.1007/s00445-016-1058-8}{10.1007/s00445-016-1058-8}
\bibitem{kozono2009} Kozono, T.\ \& Koyaguchi, T. (2009).
  \emph{Effects of relative motion between gas and liquid on 1-dimensional
  steady flow in silicic volcanic conduits}. JVGR 180:37--48.
  \href{https://doi.org/10.1016/j.jvolgeores.2008.11.006}{10.1016/j.jvolgeores.2008.11.006}
\bibitem{slezin2003} Slezin, Yu.\,B. (2003).
  \emph{The mechanism of volcanic eruptions (a steady state approach)}.
  JVGR 122:7--50.
  \href{https://doi.org/10.1016/S0377-0273(02)00464-X}{10.1016/S0377-0273(02)00464-X}
\bibitem{stuart2010} Stuart, A.\,M. (2010). \emph{Inverse problems: a Bayesian
  perspective}. Acta Numerica 19:451--559.
  \href{https://doi.org/10.1017/S0962492910000061}{10.1017/S0962492910000061}
\end{thebibliography}

\end{document}
